\section{Introduction}
\label{1-sec:1}

This chapter motivates the thesis and establishes its scope. \autoref{1-subsec:1.1 background} establishes why thermal flexibility in Norwegian district heating has become valuable to the power system, and how the current regime locks it out, \autoref{1-subsec:1.2 problem statement} states the split-incentive problem which contributes to the lockout, \autoref{1-subsec:1.3 research questions} translates it into a main research question and three sub-questions. \autoref{1-subsec:1.4 scope and limitations} draws the case and scope boundaries, \autoref{1-subsec:1.5 contribution} states the contribution claim, and \autoref{1-subsec:1.6 structure} outlines the remaining chapters.

%-----------------------------------------------------------------------
\subsection{Background and Motivation}
\label{1-subsec:1.1 background}
The marginal value of flexibility in the Norwegian power system is increasing, driven by three things. Electricity prices are projected to become more volatile, year to year, week to week and within the day. The transmission and distribution grid needs expansion, and variable renewable generation capacity continues to grow \cite{NVE_langsiktig_2023}. In the Oslo area, network operators have received applications for roughly 1,000 MW of new load connections, of which only about 200 MW have reserved capacity \cite{Thema_fjernvarme_2024}. 

District heating was non-existent in Norway before the 1980s since abundant reservoir hydropower laid the groundwork for an electricity-based heating supply, and Norway put its efforts into electricity-based heating technologies while its Nordic neighbours developed district heating \cite{sandberg_framework_2018}. That structural advantage is eroding as electricity demand outpaces available grid capacity \cite{NVE_langsiktig_2023}. The Norwegian flexibility logic is also distinct from the continental one: in Denmark and Germany, power-to-heat is valued for taking up temporary surplus generation from variable renewables \cite{bloes_power--heat_2018}, whereas additional district heating in Norway frees hydropower for export and grid balancing. In a national reference-system simulation, hydropower generation capacity is the binding constraint on exporting at the interconnection limit in 149 hours of the year, and shifting roughly a quarter of individual electric heating to district heating removes that constraint in every hour \cite{askeland_balancing_2019}. This is a technical potential at national scale rather than an Oslo-specific or market outcome.

The valuation of that flexibility sits inside three international strands that \autoref{2-sec:2} assembles. The first is the definition of energy flexibility itself, which the reviews of the field have not settled: a SINTEF synthesis finds none of sixteen definitions satisfying its three criteria \cite{degefa_comprehensive_2021}, and a systematic review finds that roughly $42\%$ of building-flexibility studies give no formal definition at all \cite{li_residential_2021}. The second is the tariff-design literature, which shows that the structure of the grid tariff decides whether power-to-heat runs: capacity-based tariffs can remove electric-boiler operation from a Norwegian plant's winter weeks \cite{kirkerud_impacts_2016}, a revenue-neutral capacity tariff would unlock a fifth more power-to-heat capacity across the Nordic system \cite{bergaentzle_cross-sector_2022}, and the competitive window of an electric boiler moves from a few per cent of hours to most of them with the tariff design in \cite{skytte_design_2017}. The third is the paired valuation of one asset from two perspectives, a Danish island system in which the operator's storage optimum falls far below what would still benefit the island system \cite{ostergaard_business_2019}, and an Austrian seasonal-storage study whose societal layer is macroeconomic rather than a welfare balance \cite{moser_socioeconomic_2018}. Across the district heating benefit literature as a whole, flexibility is monetised in none of 123 reviewed studies \cite{ahmed_district_2025}. The Norwegian case is where those strands meet a hydropower system, and no study has priced its flexibility from both perspectives at once.

District heating operators in Norway control a latent flexibility resource through electric boilers, large-scale heat pumps, and accumulator thermal storage. Both boiler fuel-switching and accumulator tanks are recognised as national sources of consumer flexibility \cite{energikommisjonen_mer_2023}, but the current regime locks that resource out on two fronts. In the reserve markets, access is administratively rationed: Statnett caps the reserve volume that can prequalify from grids below 110 kV per defined capacity area, so participation is bounded by allocation rather than by asset scale or qualification \cite{statnett_innforing_2026}. For new capacity, electric-boiler connections above 5 MW route into the grid-connection queue before any market question arises, and the operator usually installs boilers from 10 MW upward, so every addition it plans enters that queue (operator interview, Hafslund Celsio, 25 June 2026). The statutory machinery is a maturity assessment of every connection request of 1 MW or more, alongside a duty on network companies to reserve capacity for smaller consumption customers \cite{energidepartementet_forskrift_2019}. The capacity component of the grid tariff penalises increased peak electricity draw in its standard form, where the demand charge is levied on the highest measured demand within a month \cite{kirkerud_impacts_2016}, but is largely neutralised in practice for Hafslund, whose boiler fleet runs on interruptible connection terms under which the grid company may disconnect it (operator interview, Hafslund Celsio, 25 June 2026). It remains binding for firm connections and for capacity additions.


%-----------------------------------------------------------------------
\subsection{Problem Statement}
\label{1-subsec:1.2 problem statement}
Societal-scale analysis shows that flexibility in district energy and electricity systems benefits society. That benefit is not reflected in industry returns, which produces a split-incentive problem that should be addressed \cite{sneum_barriers_2021}. Whether that condition holds for Norwegian district heating is a question this thesis takes up. District heating (DH) operators who invest in, and operate, flexible assets like electric boilers, thermal storage, and heat pumps, bear the full financial cost of supplying that flexibility: capital expenditure, financing constraints, the operational complexity of market participation, and grid capacity charges \cite{sneum_barriers_2021}. The potential benefits of activating that flexibility, however, accrue elsewhere: to distribution system operators through avoided grid reinforcement, to transmission system operators through freed balancing capacity, and to consumers through potential lower system prices during peak periods. No functional transfer mechanism currently exists to redirect a share of those societal gains back to the DH operators, who carry the financial burden \cite{Thema_fjernvarme_2024}.

The consequence is a structurally sub-optimal equilibrium, and the divergence it produces takes two forms that call for different instruments. The first is that the value flexibility creates primarily reaches parties who did not pay for it, which is a question of how a total that exists is divided. The same conditions can also stop value from being created at all, by withholding the market access that flexible operation depends on, which is a question of what the total is allowed to be. The incentive operators face to offer flexibility accordingly falls short of what the power system would gain from it. Where they hold no reserve access at all, the allocation withholds the revenue that would create that incentive. Closing this divide requires understanding both sides of the ledger: the commercial costs and revenues that determine operator behaviour, and the socio-economic value that justifies public interest in changing those conditions. This thesis examines both, measures the two forms of the divergence separately, and analyses what regulatory and market-design changes would be required to align them.


%-----------------------------------------------------------------------
\subsection{Objectives and Research Questions}
\label{1-subsec:1.3 research questions}

\paragraph*{Main Research Question:} 
To what extent do current market and regulatory conditions produce a divergence between the commercial and socio-economic value of thermal flexibility in Norwegian district heating, and what mechanisms could reduce that divergence?

\paragraph*{Sub-Questions}
\begin{enumerate}
    \item What are the commercial costs and revenues associated with operating thermal flexibility assets in the Oslo district heating system under current market and tariff conditions?
    \item What is the socio-economic value of the same flexibility, and how robust is the triangulated estimate?
    \item What specific regulatory, market-design, and grid-tariff features contribute the most to the divergence between the two perspectives, and which are most amenable to change?
\end{enumerate}


%-----------------------------------------------------------------------
\subsection{Scope and Limitations}
\label{1-subsec:1.4 scope and limitations}
The thesis takes Oslo's district heating system, operated by Hafslund Celsio, as the primary single case. Trondheim and Bergen are retained as comparative system context only. The flexibility assets in scope are electric boilers, heat pumps and accumulator thermal storage. Small, local heating systems, end-user price elasticity, demand-side response, and detailed technical asset design are out of scope, as is an independent socio-economic cost-benefit analysis. The socio-economic valuation is instead drawn from existing literature and triangulated across those sources rather than constructed anew \cite{ahmed_district_2025}. Throughout the analysis, the district heating operator is treated as a price taker in the power market and, by extension, in the reserve market. Its dispatch responds to prices but does not influence them. The extension to the reserve market is materially weaker than the power-market assumption, and it is what makes the highest reserve-access results an upper bound rather than an estimate. The assumption is defended in \autoref{3-subsec:3.4 operator model} and matches the operator's own description of its dispatch practice (operator interview, Hafslund Celsio, 25 June 2026).

The reduction to a single primary case is defended in \autoref{3-subsec:3.1 research design}. \autoref{3-subsec:3.10 limitations} sets out the limitations the design carries by construction. The remaining scope choices, and the work they leave for others, are taken up in \autoref{5-subsec:5.5 future work}.

%-----------------------------------------------------------------------
\subsection{Contribution}
\label{1-subsec:1.5 contribution}
The thesis offers a paired commercial and socio-economic valuation of thermal flexibility for a single Norwegian district heating system, quantifying the gap between the two perspectives rather than presenting them as parallel analyses. It separates the value the operator creates but does not capture from the value the current access allocation prevents it from creating. The operator-side estimate is produced by a deterministic Python and PuLP dispatch model run against current market and tariff conditions. The system-side estimate triangulates published socio-economic figures across multiple sources. The closest precedent applies a comparable paired structure to a Danish biomass-based island system \cite{ostergaard_business_2019}. A second study pairs a plant-level monetisation with a regional macroeconomic layer for an Austrian seasonal-storage case, but that layer reports gross regional product, net export balance and employment effects rather than a welfare balance, so it transfers as a method rather than as a paired comparator \cite{moser_socioeconomic_2018}. Both are methodologically reusable but not substantively comparable to the hydropower-coupled Norwegian context. To the author's knowledge, a paired commercial and socio-economic valuation has not previously been applied to a Norwegian district heating system at the operator level with explicit quantification of the gap.

%-----------------------------------------------------------------------
\subsection{Thesis Structure}
\label{1-subsec:1.6 structure}

\autoref{1-sec:1} states the problem, the research questions and the scope, and locates the contribution.

\autoref{2-sec:2} establishes the theoretical framework: energy-system flexibility concepts, district heating as a flexibility resource, the barriers to using it, the Norwegian context, the conceptual definitions of commercial and socio-economic value, and a literature review that maps the field and locates this thesis within it.

\autoref{3-sec:3} describes the methodology: the case-study design and the evidence it rests on, the technical flexibility envelope that sizes the physical resource before any price is applied to it, the optimisation model, the counterfactual and scenario grid the model is run across, the socio-economic valuation triangulation, and the gap-analysis framework.

\autoref{4-sec:4 results} reports the results in the order the method produces them. It covers how the dispatch uses the physical resource and how the model fares against the public record and its own tests. It then reports the operator's ledger across the capability ladder and the scenario grid, and the socio-economic valuation of the same dispatch. It divides the socio-economic total between the operator and the parties that receive the remainder, with the tariff and access barriers priced and attributed across the barrier types. It closes on the sensitivity of every result to its inputs and to the delivery year.

\autoref{5-sec:5 discussion} interprets those results against the barriers of \autoref{2-subsec:2.3 barriers}, draws out the policy and market-design implications and the conditions under which the findings transfer, and closes on what the study's limitations cost its conclusions and the work that follows from them.

\autoref{6-sec:6 conclusion} answers the research questions and states what the thesis establishes and what it does not.

\newpage